\section{Parte B: Circuito RL de Corriente Alterna}

\subsection{Objetivos y Fundamento Teórico}
El propósito de esta segunda etapa es determinar experimentalmente la resistencia interna ($r$) y la inductancia ($L$) de una bobina real en un circuito serie de corriente alterna alimentado a frecuencia industrial ($50\text{ Hz}$). 

A diferencia de un inductor ideal, una bobina real posee una resistencia óhmica asociada debido al bobinado de cobre. Por lo tanto, la impedancia del circuito serie $R-L-r$ se expresa como:
\begin{equation}
Z = \sqrt{(R+r)^2 + (\omega L)^2}
\end{equation}

El desfase angular general del circuito viene dado por:
\begin{equation}
\tan(\varphi) = \frac{\omega L}{R+r}
\end{equation}

\subsection{Resultados Experimentales y Análisis}
Utilizando una fuente de tensión alterna de $12\text{ V}$ a $50\text{ Hz}$ y un resistor nominal de aproximadamente $50\ \Omega$, se midieron de forma directa la corriente eficaz del circuito y las caídas de tensión en cada componente mediante instrumentación de laboratorio.

Los valores experimentales obtenidos y las magnitudes calculadas mediante la resolución del diagrama fasorial se detallan a continuación:

\begin{table}[h!]
\centering
\caption{Mediciones y parámetros determinados para el circuito RL real.}
\label{tab:parteB}
\begin{tabular}{cccccc}
\toprule
$V_R$ (V) & $V_{zL}$ (V) & $V_{total}$ (V) & $I$ (mA) & $r$ ($\Omega$) & $L$ (mH) \\
\midrule
5.36 & 9.08 & 12.95 & 105.6 & 49.93 & 222.817 \\
\bottomrule
\end{tabular}
\end{table}

A través del método del paralelogramo para las tensiones fasoriales, se dedujo un ángulo de fase $\alpha = 34.806^{\circ}$. Esto permitió proyectar la caída de tensión de la bobina en sus componentes intrínsecas: una caída resistiva pura $u_r \approx 5.273\text{ V}$ y una caída inductiva pura $V_L \approx 7.392\text{ V}$, logrando aislar con éxito los parámetros reales del componente.